\begin{figure}[htbp]
\centering
\begin{tikzpicture}[>=Stealth, line width=0.8pt, font=\small,
nodo/.style={circle, fill=white, draw=black,
	inner sep=0pt, minimum size=5pt, line width=0.8pt},
cint/.style={circle, fill=white, draw=black,
	inner sep=0pt, minimum size=6pt, line width=0.8pt},
fantasma/.style={thin, gray, dashed},
tacca/.style={thin, gray!70}]

\def\a{1.0}
\def\tw{0.22}\def\thA{0.38}\def\thR{0.32}
\def\bw{0.16}\def\gw{0.30}\def\sr{0.08}

%% riferimento indeformato, tratteggiato
\draw[fantasma] (0,0) -- ({5*\a},0);

%% configurazione del meccanismo
\draw[line width=2.2pt] (0,0) -- ({3*\a},-0.9);
\draw[line width=2.2pt] ({3*\a},-0.9) -- ({5*\a},0);

\node[above=2pt] at (0,0)             {$A$};
%\node[below=3pt] at ({2*\a},-0.6)     {$B$};
\node[below=3pt] at ({3*\a},-0.9)     {$C$};
\node[above=2pt] at ({5*\a},0)        {$D$};

%\node at ({1*\a},-0.15)  {$\mathcal{C}_1$};
%\node at ({4*\a},-0.6)   {$\mathcal{C}_2$};

%% angolo virtuale theta_1 in A
\draw[tacca] (0,0) -- (0.9,0);
\draw[thin] (0.75,0) arc (0:-16.7:0.75);
\node[below right=1pt] at (0.9,.09) {$\vartheta_1$};

%% angolo virtuale theta_2 in C, contro una tacca locale orizzontale
\draw[tacca] ({3*\a-0.4},-0.9) -- ({3*\a+0.4},-0.9);
\draw[thin] ({3*\a+0.35},-0.9) arc (0:24.2:0.35);
\node[above right=1pt] at ({3*\a+0.4},-1.1) {$\vartheta_2$};

%% cerniera in A, invariata
\draw (0,0) -- ++(-\tw,-\thA) -- ++(2*\tw,0) -- cycle;
\draw (-\gw,-\thA) -- (\gw,-\thA);
\foreach \x in {-0.24,-0.12,0,0.12,0.24}
\draw[thin] (\x,-\thA) -- ++(-0.09,-0.11);
\node[nodo] at (0,0) {};

%% carrello in B soppresso: fantasma nella posizione originaria
\draw[fantasma] ({2*\a},0) -- ++(-\tw,-\thR) -- ++(2*\tw,0) -- cycle;
\node[gray, above=2pt] at ({2*\a},0) {\footnotesize(soppresso)};
%\node[nodo] at ({2*\a},-0.6) {};

%% cerniera interna, nella posizione attuale
\node[cint] at ({3*\a},-0.9) {};

%% carrello in D, invariato (coincide con la posizione di riferimento)
\draw ({5*\a},0) -- ++(-\tw,-\thR) -- ++(2*\tw,0) -- cycle;
\draw ({5*\a-\gw},-\thR) -- ({5*\a+\gw},-\thR);
\draw ({5*\a-\bw},-\thR-0.12) circle (2.4pt);
\draw ({5*\a+\bw},-\thR-0.12) circle (2.4pt);
\draw ({5*\a-\gw},-\thR-0.24) -- ({5*\a+\gw},-\thR-0.24);
\foreach \x in {-0.24,-0.12,0,0.12,0.24}
\draw[thin] ({5*\a+\x},-\thR-0.24) -- ++(-0.09,-0.11);
\node[nodo] at ({5*\a},0) {};

%% reazione incognita Y_B, applicata nel punto attuale
\draw[-{Latex[length=2.8mm]}, thick] ({2*\a},-1.2) -- ({2*\a},-0.6);
\node[below=2pt] at ({1.7*\a},-0.85) {$Y_B$};

%% carichi attivi, applicati nei punti attuali
\draw[-{Latex[length=2.8mm]}, thick] ({1*\a},0.65) -- ({1*\a},0.1);
\node[right=2pt] at ({0.2*\a},0.5) {$2F$};

\draw[-{Latex[length=2.8mm]}, thick] ({4*\a},0.6) -- ({4*\a},0.1);
\node[right=2pt] at ({4*\a},0.5) {$F$};

\end{tikzpicture}
\caption{Il meccanismo usato per determinare $Y_B$: si sopprime il
carrello in $B$. $\mathcal{C}_1$ ruota di $\vartheta_1$ attorno alla
cerniera in $A$; $\mathcal{C}_2$ accompagna il moto ruotando di
$\vartheta_2=-\tfrac32\vartheta_1$ attorno a $C$, in modo da lasciare
fermo il punto $D$, dove il carrello non è stato toccato.}
\label{fig:gerber_mecc_B}
\end{figure}
